\documentclass{report}
\usepackage{amsfonts, amsmath}
\newcommand\R{{\mathbb R}}
\newcommand\Q{{\mathbb Q}}
\newcommand\N{{\mathbb N}}
\newcommand\ve{\varepsilon}
\newcommand\Co{{\mathcal C}}
\pagestyle{empty}
\begin{document}
\large
\centerline{\Huge Fisica Matematica II}
\renewcommand{\thefootnote}{ } 
\footnote{\sl Avete 2:30 ore di tempo. Ogni esercizio
vale dieci punti. La sufficienza si ottiene con un punteggio $\geq$
18. Solo le risposte chiaramente giustificate saranno prese in
considerazione.}
\renewcommand{\thefootnote}{\arabic{footnote}} 
\vskip.1cm
\centerline{\Large Appello 06-10-2003}
\vskip1cm
\begin{enumerate}
\item Si determini una soluzione dell'equazione
\[
\begin{split}
&yu_x-xu_y=u^2\\
&u(x,0)=x^2\quad x\in\R
\end{split}
\]
e se ne discuta dominio e unicit\`a.
\item Si trovi una soluzione dell'equazione
\[
\begin{split}
&\Delta u(x,y)=x^2y^2\quad \quad x\in(0,1)\,\;y\in(0,2),\\
&u(x,y)=0\quad \quad\forall\; (x,y)\in\partial([0,1]\times[0,2]).
\end{split}
\]
Se ne discuta unicit\`a e regolarit\`a.
\item Si consideri l'equazione
\[
\begin{split}
&u_{tt}=u_{xx} \quad (t,x)\in\R_+\times [0,2]\\
&u(x,0)=x;\; \;\;u_t(x,0)=0\\
&u(0,t)=0;\; \;\;u(2,t)=1.
\end{split}
\]
Si determini la soluzione e si mostri che l'energia si conserva.
\item Sia $u$ la soluzione dell'equazione del calore
\[
\begin{split}
&u_t(x,t)=u_{xx}(x,t)\quad (x,t)\in\R\times\R_+\\
&u(x,0)=\frac{x^2}{x^2+1}.
\end{split}
\]
Si calcoli,
\[
\lim_{t\to\infty}u(x,t).
\]

\end{enumerate}
\end{document}
